\chapter{Implementation}

The QuantumHop system is implemented using Python for the backend and React for the frontend. The backend is structured as a Python package (\texttt{backend/}) with separate modules for each concern. The frontend is a Vite-based React application served on port 5173, which proxies all \texttt{/api} requests to the Flask server on port 5000. All inter-machine communication is handled exclusively through raw TCP sockets on port 5010; the Flask API is used only for browser-to-backend communication on the same machine.

\section{Implementation Environment}

\begin{enumerate}
    \item \textbf{Python 3.10+} with \texttt{flask}, \texttt{flask-cors}, and \texttt{cryptography} libraries.
    \item \textbf{Node.js 18+} with \texttt{npm} for the Vite/React frontend.
    \item \textbf{Windows PowerShell} for the launcher batch script.
    \item \textbf{Firewall rules} to allow inbound TCP traffic on port 5010 and UDP traffic on port 5555.
\end{enumerate}

\section{Module Descriptions}

\subsection{bb84.py --- Quantum Key Distribution Engine}

The \texttt{establish\_key()} function is the core of the system. It generates \texttt{rounds = max(KEY\_BITS * 4, SAMPLE\_SIZE * 4)} random Alice bits and bases, and an independent set of Bob bases. When \texttt{eavesdrop=True}, a third Eve basis set is generated; Eve's measurement randomises bits where her basis differs from Alice's, and Bob further randomises bits where his basis differs from Eve's. This double randomisation produces the characteristic $\approx$25\% error rate signature.

Basis sifting keeps only the positions where Alice and Bob chose the same basis. A configurable sample of sifted bits is compared to compute the error rate. The remaining sifted bits are hashed with SHA-256 to produce the 32-byte AES key. The function returns a \texttt{BB84Result} dataclass containing all previews needed by the dashboard.

\begin{verbatim}
@dataclass(frozen=True)
class BB84Result:
    key: bytes
    error_rate: float
    accepted: bool
    sifted_bits: list[int]
    compared_bits: int
    generated_bits: int
    matching_bases: int
    alice_basis_preview: str
    bob_basis_preview: str
    keep_preview: str
\end{verbatim}

\subsection{crypto\_utils.py --- AES-256-CBC Encryption}

The \texttt{encrypt\_message()} function accepts a plaintext string and a key (bytes or hex string), derives a 32-byte AES key via SHA-256, generates a random 16-byte IV, applies PKCS7 padding, and encrypts using AES-CBC. The result is a dictionary with Base64-encoded \texttt{iv} and \texttt{ciphertext} fields suitable for JSON serialisation.

The \texttt{decrypt\_message()} function reverses this process. It accepts the payload dictionary and the key material, derives the AES key, and decrypts and unpads the ciphertext. A critical design decision is that the sender's packet carries the derived key in hex so that the receiver can decrypt with the same key without re-running BB84, since a simulated BB84 on two machines would produce different keys.

\subsection{peer\_socket.py --- Raw TCP Socket Layer}

This module implements both the outbound sender and the inbound server.

\textbf{Sender side}: \texttt{send\_message\_to\_peer()} calls \texttt{bb84.establish\_key()}, encrypts the message, constructs the JSON envelope, and dispatches it via \texttt{\_send\_envelope()} which opens a \texttt{socket.create\_connection()} to the target IP and port, calls \texttt{sendall()}, signals end-of-send with \texttt{shutdown(SHUT\_WR)}, then reads the response in a loop.

\textbf{Server side}: \texttt{run\_peer\_server()} creates a listening TCP socket bound to \texttt{0.0.0.0:5010}. For each accepted connection, it spawns a daemon thread running \texttt{\_handle\_client()}. This function reads the full envelope (loop until empty recv), parses the JSON, routes to either \texttt{\_handle\_message()} or \texttt{\_handle\_relay()}, wraps the result, and sends it back before closing the connection.

\textbf{MITM relay}: When a relay envelope arrives, \texttt{\_handle\_relay()} logs the interception event and forwards a modified envelope (with \texttt{attackMode="mitm"}) to the real target using another \texttt{\_send\_envelope()} call.

\subsection{peer\_discovery.py --- UDP Broadcast Discovery}

Two daemon threads run continuously. The broadcast thread sends a JSON payload containing the machine name, local IP, Flask port, and socket port to the LAN broadcast address on port 5555 every 3 seconds. The listen thread binds to port 5555 and updates an in-memory peer registry keyed by IP. Peers not heard from within 12 seconds are evicted. The \texttt{get\_peers()} function returns the current registry excluding self.

\subsection{attack\_detector.py --- Nonce Registry and Error Rate Check}

The \texttt{register\_nonce()} function inserts a nonce into a thread-safe dictionary with the current timestamp. If the nonce already exists and has not expired (TTL = 60 seconds), it returns a failure result. The \texttt{is\_error\_rate\_attack()} function compares the error rate against \texttt{ERROR\_THRESHOLD} (0.15).

\subsection{app.py --- Flask API Gateway}

Flask serves the API endpoints that the React dashboard calls:
\begin{itemize}
    \item \texttt{POST /api/send-to-peer} --- Validates the request, calls \texttt{send\_message\_to\_peer()}, returns the result JSON.
    \item \texttt{GET /api/peers} --- Returns the current peer list from \texttt{peer\_discovery.get\_peers()}, plus a self-entry.
    \item \texttt{GET /api/inbox} --- Returns all messages from \texttt{inbox.get\_messages()}.
    \item \texttt{GET /api/events} --- Returns all logged events from \texttt{logger.get\_events()}.
    \item \texttt{POST /api/reset} --- Clears inbox, events, and nonce registry.
\end{itemize}

\subsection{React Frontend}

The frontend is built with React 18 and Vite. The \texttt{App.jsx} component manages state for the active tab, peer list, inbox, events, selected peer, attack mode, relay IP, and send result. It polls all four API endpoints every 2 seconds.

The \textbf{Real Network} tab renders a message composition panel (target dropdown, mode dropdown, message input, send button) and two panels below: the \textbf{Network Peers} list showing each discovered peer with name, IP, and online indicator, and the \textbf{Inbox} panel showing each received message as a card with sender, content or block reason, BB84 error rate badge, and a coloured attack alert banner for blocked messages.

\section{Launcher Script}

The \texttt{runeverything.bat} launcher script performs the following steps in sequence:

\begin{enumerate}
    \item Scans ports 5000, 5010, 5173, and 5555 using \texttt{netstat} and forcefully kills any occupying process with \texttt{taskkill /F /PID}.
    \item Activates the Python virtual environment and runs \texttt{pip install -r requirements.txt}.
    \item Starts the Flask backend in a new terminal window using \texttt{python -m backend.app}.
    \item Waits 2 seconds, then starts the Vite frontend in another terminal window using \texttt{npm run dev}.
\end{enumerate}

This ensures a clean, reproducible startup every time, regardless of whether a previous session was terminated cleanly.

\section{Challenges and Solutions}

\begin{enumerate}
    \item \textbf{Silent server crash:} The original \texttt{\_handle\_client()} had no try-except around the handler call. An unhandled \texttt{TypeError} (duplicate keyword argument in a logger call) caused the server thread to die silently, returning empty bytes to the client and producing a JSON parse error. This was fixed by wrapping the handler dispatch in a try-except that always sends a valid JSON error response.

    \item \textbf{Double SHA-256 hashing:} The \texttt{derive\_aes\_key()} function re-hashes its input with SHA-256. Since the BB84 key is already a 32-byte SHA-256 digest, passing it as a hex string caused double-hashing, making decryption fail. The fix was to pass the key as raw bytes using \texttt{bytes.fromhex()} before calling \texttt{decrypt\_message()}.

    \item \textbf{Duplicate keyword argument:} A \texttt{logger.emit\_event()} call passed \texttt{errorRate=error\_rate} both explicitly and via \texttt{**bb84\_details}, causing a \texttt{TypeError}. The solution was to pass \texttt{bb84\_details} as a single keyword argument rather than unpacking it.

    \item \textbf{Port hardcoding:} The frontend originally called \texttt{sendToPeer()} with port 5000 (the Flask port) instead of 5010 (the socket port), causing Flask to receive a raw TCP request it could not parse. The fix was to resolve the socket port dynamically from the peer discovery object.
\end{enumerate}
